\documentclass{article}

\usepackage{PRIMEarxiv}

\usepackage[utf8]{inputenc}
\usepackage[T1]{fontenc}
\usepackage{hyperref}
\usepackage{url}
\usepackage{booktabs}
\usepackage{array}
\usepackage{amsfonts}
\usepackage{amsmath}
\usepackage{nicefrac}
\usepackage{microtype}
\usepackage{fancyhdr}
\usepackage{graphicx}
\usepackage{xcolor}
\usepackage{tikz}
\usetikzlibrary{arrows.meta,positioning,shapes.geometric,fit,calc,decorations.pathreplacing,backgrounds}
\usepackage{enumitem}
\usepackage{algorithm}
\usepackage{algpseudocode}
\usepackage{pifont}
\usepackage{amssymb}
\usepackage{xspace}
\usepackage[text=PREPRINT,color=gray!25,scale=3]{draftwatermark}
\graphicspath{{media/}}

\newcommand{\cmark}{\ding{51}}
\newcommand{\xmark}{\ding{55}}
\newif\ifanonymous
% \anonymoustrue
\anonymousfalse
% shared/macros.tex — Identity-masked macros for both arXiv and symposium versions
% Toggle anonymization: set \anonymoustrue before \input{macros} for double-blind review,
%                       set \anonymousfalse before \input{macros} for arXiv

\ifanonymous
  \newcommand{\toolname}{{\sc SMBGrain}\xspace}
  \newcommand{\anonrepo}{\url{https://anonymous.4open.science/r/XXXX}}
  \newcommand{\authorblock}{%
    \author{Anonymous Submission \\ Paper ID: TBD}%
  }
\else
  \newcommand{\toolname}{{\sc ksmbdzzer}\xspace}
  \newcommand{\anonrepo}{\url{https://github.com/yskzalloc/kcov-dataflow}}
  \newcommand{\authorblock}{%
    \author{Yunseong Kim \\
      Ericsson Software Technology \\
      \texttt{yunseong.kim@est.tech}}%
  }
\fi

% Technical shorthands
\newcommand{\kcovdf}{{\sc kcov-dataflow}\xspace}
\newcommand{\ksmbd}{\texttt{ksmbd}\xspace}
\newcommand{\sancov}{\texttt{SanitizerCoverage}\xspace}
\newcommand{\tracepc}{\texttt{trace-pc}\xspace}
\newcommand{\traceargs}{\texttt{trace-args}\xspace}
\newcommand{\traceret}{\texttt{trace-ret}\xspace}
\newcommand{\tracecmp}{\texttt{trace-cmp}\xspace}
\newcommand{\kcov}{\texttt{KCOV}\xspace}
\newcommand{\kasan}{\texttt{KASAN}\xspace}
\newcommand{\syzkaller}{\texttt{syzkaller}\xspace}
\newcommand{\libfuzzer}{\texttt{libFuzzer}\xspace}
\newcommand{\virtme}{\texttt{virtme-ng}\xspace}

% Data-flow tuple notation
\newcommand{\dftuple}[4]{\ensuremath{\langle #1, #2, #3, #4 \rangle}}
\newcommand{\callerid}{\ensuremath{\mathit{PC}_{\mathrm{caller}}}}
\newcommand{\calleeid}{\ensuremath{\mathit{PC}_{\mathrm{callee}}}}

% Results that are not yet measured on the paper's toolchain. Every \pending{}
% must be replaced by a measured value before submission; grep for it.
\newcommand{\pending}[1]{\textcolor{red}{\textbf{[#1]}}}

% Listing style
\usepackage{listings}
\lstset{
  basicstyle=\ttfamily\footnotesize,
  breaklines=true,
  frame=single,
  numbers=left,
  numberstyle=\tiny,
  captionpos=b,
}


%% ---------------------------------------------------------------------------
%% STATUS OF THE EVALUATION (read before submitting)
%%
%% Every number printed in red via \pending{} is NOT yet measured on the
%% paper's toolchain. The engine-comparison campaigns recorded so far
%% (~/ksmbdzzer-logs/2026-08-28_*, 2026-09-01_*, 2026-09-02_*) are all
%% INVALID: their summaries report kern_pcs = 0 for every arm, and the logs
%% show "REMOTE_ENABLE(0x4b440001) failed ... Bad address" (Section~\ref{sec:correction}).
%% Measured facts used in this draft: the 2026-07-22 grain selftest
%% (208 WORKS / 0 DEAD / 0 BAIL / 2 EXCL), the upstream fixes and CVEs, the
%% source-code constants, and the kcov-dataflow cost figures taken from the
%% foundational paper.
%% ---------------------------------------------------------------------------

% Header
\pagestyle{fancy}
\thispagestyle{empty}
\renewcommand{\shorttitle}{Dataflow-Steered SMB-Procedure Fuzzing of the Linux In-Kernel SMB Server}
\renewcommand{\headeright}{A Preprint}

%% Title
\title{Deep by Construction: Dataflow-Steered, Procedure-Level Fuzzing of the Linux In-Kernel SMB Server}

\authorblock
\date{}

\begin{document}
\maketitle

\begin{abstract}
The Linux in-kernel SMB3 server, \texttt{ksmbd}, runs a network-reachable
protocol stack in ring~0. Fuzzing it means clearing three obstacles. A
\emph{protocol wall}: file-system logic is reached only after a valid
negotiate, NTLMv2 session setup, tree connect and \texttt{CREATE}. A
\emph{value wall}: once the handlers are covered, edge coverage cannot tell a
benign length from a malicious one. An \emph{oracle wall}: the most damaging
server defect, modifying a file the client may not touch, does not crash. I
present \toolname, a fuzzer that works on SMB \emph{procedures} rather than
system calls. A \emph{grain} is a valid scenario whose authenticated prefix
is never mutated; only its final protocol data unit is fuzzed, so every input
starts past the parser. 210 grains cover all 19 SMB2 commands, every FSCTL
and information class, the \texttt{CREATE} contexts, SMB3 transforms,
SMBDirect and SMB1. Feedback comes from \kcovdf, a compiler-based KCOV
extension that records function arguments, return values and comparison
operands. A handle derived from each connection's loopback address routes
those records to that connection's harness. The harness folds them into
\libfuzzer's counters as saturating path features plus bounded, first-seen
value features, and feeds the comparison operands to an input-to-state
mutator. An effect-based contract oracle checks write-side integrity against
a root-owned victim, and a stability oracle flags silent state corruption.
208 of the 210 grains reach \texttt{ksmbd} code. The campaigns found six
vulnerabilities, all assigned CVEs and fixed upstream. The aim is not to beat
edge-coverage fuzzers on their own metric, but to guide mutation with what the
kernel's functions actually receive and return.
\end{abstract}

\keywords{Kernel Fuzzing \and SMB \and ksmbd \and Data-Flow Feedback \and Compiler Instrumentation \and Input-to-State Correspondence \and Logic-Bug Oracles}


\section{Introduction}\label{sec:intro}

\ksmbd~\cite{ksmbddoc} is the SMB3 file server built into Linux. It parses
untrusted network input in kernel context, and its request handlers run on
kernel workqueues with the privileges needed to act on any file in an exported
tree. A single defect there can be reached remotely, with the attacker sending
the bytes and the kernel parsing them. The server's history shows the
consequence: CVE-2022-47939~\cite{cve202247939}, a use-after-free reachable
during tree disconnect, was rated critical. Since then, file-table
lifetimes, IPC message validation and the SMBDirect transport have all
produced memory-safety bugs. The security-critical paths are well known; the
open question is how to reach them systematically and judge what happens
there.

A fuzzer faces three distinct obstacles in this setting.

\paragraph{The protocol wall.} \ksmbd does nothing interesting with a request
until a session exists. The client must negotiate a dialect, complete a
two-leg NTLMSSP exchange whose response is an HMAC over a server
challenge~\cite{msnlmp}, connect to a share, and open a handle. Signing or
encryption may be required on top. Random bytes, and mutations of a captured
trace that break the session state, are rejected in
\texttt{ksmbd\_smb2\_check\_message()} or in session lookup, long before any
file-system code runs. Syscall-level kernel fuzzers such as
\syzkaller~\cite{vyukov2015syzkaller} can open a TCP socket to port~445, but
they would have to generate the NTLMv2 proof and keep the message identifiers,
credits and session identifiers consistent across calls.

\paragraph{The value wall.} When valid requests do reach the handlers, edge
coverage soon runs out of new edges. \texttt{smb2\_write()} with a length of
eight and with a length of \texttt{0xFFFFFFFF} take the same edges up to the
bounds check. A \texttt{CREATE} asking for \texttt{FILE\_READ\_DATA} and one
asking for \texttt{FILE\_WRITE\_DATA|DELETE\_ON\_CLOSE} also look the same to
\tracepc. These argument values decide what the server does, and edge coverage
does not record them~\cite{kim2026kcovdataflow}.

\paragraph{The oracle wall.} The defect class that matters most for a file
server is an \emph{integrity violation}: the server truncates, overwrites or
deletes a file the authenticated user has no right to modify. Such a bug
returns \texttt{STATUS\_SUCCESS} and leaves no \kasan report or lockdep splat.
Crash oracles miss it, and a naive ``request succeeded although access was
not granted'' check produces false positives, because \ksmbd authorizes at the
POSIX layer and an owner may legitimately operate on its own file
(Section~\ref{sec:correction}).

\medskip
\toolname is designed around these three obstacles. It fuzzes at the
granularity of \emph{SMB procedures}, steers with \emph{kernel values}, and
judges results with \emph{effect-based contract oracles}. It operates through
three layers:
\begin{enumerate}[nosep]
  \item \textbf{Grains (Section~\ref{sec:grain}).} A grain is a working SMB
        scenario whose authenticated prefix is fixed. \libfuzzer mutates only
        the final target PDU or its parameters, so every input reaches deep
        handler code by construction. 210 grains, ordered by kernel data-path
        depth, cover the whole SMB1/2/3 procedure surface.
  \item \textbf{Value-aware kernel feedback
        (Sections~\ref{sec:route}--\ref{sec:i2s}).} \kcovdf records reach the
        right harness because each connection's records are routed by a
        handle derived from its loopback destination address. The harness
        folds them into \libfuzzer's extra counters as a path region, which
        saturates, and a value region, which adds a feature only for a
        never-before-seen value at a call site. A RedQueen-style
        mutator~\cite{aschermann2019redqueen} consumes the comparison
        operands.
  \item \textbf{Orchestration and oracles
        (Sections~\ref{sec:combine}--\ref{sec:loop}).} A four-phase round loop
        (enumerate, saturate, combine, save) runs the grain fleet in parallel,
        combines grains pairwise on a shared handle, and carries corpora
        forward from round to round. Crash, stability and contract oracles run
        as checks alongside fuzzing rather than as separate phases.
\end{enumerate}

\paragraph{Scope: a different axis, not a faster edge fuzzer.} This work does
not set out to beat existing fuzzers at their own game, and it is not a
comparison against them. Edge-coverage engines such as \syzkaller decide what
to keep using only \emph{which} code ran, and they mutate only bytes that the
fuzzer itself generated. \toolname explores a different question: what
happens when mutation is guided by knowledge of the target, namely the values
its functions actually \emph{received} and \emph{returned}? A head-to-head
benchmark against an edge-coverage engine cannot answer that question
cleanly, because the two systems differ in interface, input model, feedback
and oracle at once. A syscall-level fuzzer and a procedure-level fuzzer do not
even run the same inputs. I therefore measure the value signal \emph{inside}
one system, where everything else is held fixed (Section~\ref{sec:eval}), and
treat edge-coverage fuzzers as complementary. The direction is promising
because it opens a channel that edge coverage does not have. The kernel tells
the fuzzer which lengths, flags, handles and comparison operands its code
actually processed, and the fuzzer can push exactly those values towards
boundaries or reuse them in later requests, instead of guessing blindly. This
paper is a first, end-to-end step in that direction on a production kernel
subsystem, and the direction deserves further development.

My contributions are:
\begin{itemize}[nosep]
  \item \textbf{Procedure-level, deep-by-construction fuzzing of an in-kernel
        network server.} The grain abstraction fixes a valid authenticated
        prefix and fuzzes only the last procedure. A 210-grain fleet covers
        all 19 SMB2 commands, all 16 FSCTLs handled by \texttt{smb2\_ioctl()},
        the \texttt{SET\_INFO}/\texttt{QUERY\_INFO} classes, 11
        \texttt{CREATE} contexts, the SMB3 transforms, SMBDirect, and the SMB1
        compatibility opcodes. A coverage-measured self-test shows that 208 of
        the 210 grains reach \ksmbd handlers.
  \item \textbf{Per-connection routing of kernel dataflow.} An IP-derived
        remote handle (\texttt{0x4b440000 | low16(daddr)}) is set in
        \ksmbd's TCP transport and published from user space through
        \texttt{KCOV\_DF\_REMOTE\_ENABLE}. Each concurrent fuzzing worker
        therefore receives exactly the argument, return and comparison records
        produced while \ksmbd services \emph{its} connection, even though that
        work runs on shared kworkers.
  \item \textbf{A bounded value-coverage fold and a controlled engine
        ablation.} The fold keeps edge coverage's clean saturation signal and
        adds a first-seen value region with value-class normalization. Five
        engine arms run the same grains on the same kernel, so the
        contribution of value coverage and of input-to-state mutation can be
        \emph{measured} rather than asserted.
  \item \textbf{Effect-based integrity oracles for a file server.} A
        cross-user contract oracle checks a root-owned victim through the
        backing file system. It is accompanied by lease-grant, race-integrity,
        teardown-lifetime and fragmented-framing oracles, and by a stability
        oracle that flags a working scenario that stops working.
  \item \textbf{Upstream impact.} Six vulnerabilities were assigned CVEs and
        fixed in Linux (Table~\ref{tab:bugs}), plus one further SMBDirect
        locking fix.
\end{itemize}


\section{Background \& Motivation}\label{sec:background}

\subsection{ksmbd and its Attack Surface}

\ksmbd splits the server between the kernel and a user-space daemon,
\texttt{ksmbd.mountd}. The daemon owns the user database and the share
configuration, and answers login, tree-connect and RPC requests over a generic
netlink channel (\texttt{transport\_ipc.c}). The kernel owns everything else.
A per-connection handler thread (\texttt{ksmbd\_conn\_handler\_loop()}) reads
NetBIOS-framed PDUs from TCP, or from SMBDirect over RDMA, and queues each
request to a kworker, where \texttt{handle\_ksmbd\_work()} validates and
dispatches it to one of the 19 SMB2 command handlers in
\texttt{smb2pdu.c}. Compound requests chain several commands in one frame.
SMB3 adds a transform header for encryption and a compression transform, and
SMBDirect adds its own credit-based framing on queue pairs.

Two properties of this design shape everything that follows. First, the code
that matters runs \emph{off the client's task}: it runs on the handler thread
and on kworkers, not in the context of the process that sent the bytes.
Per-task coverage collection, the mode most fuzzers use, sees nothing; the
coverage has to be routed remotely. Second, \ksmbd authorizes in two places:
it checks the granted access mask in its SMB layer, and it then calls the VFS
with the mapped user's credentials, so POSIX permissions remain the last line
of defence. An oracle has to distinguish a genuine bypass from the case where
the SMB layer misses a check that POSIX enforces anyway.

\subsection{KCOV and kcov-dataflow}

KCOV~\cite{vyukov2016kcov} records the program counters a task executes into a
buffer mapped into user space. Its remote mode lets a kernel thread annotated
with \texttt{kcov\_remote\_start\_common(handle)} attribute its coverage to
whichever user-space collector registered that handle. \kcovdf~\cite{kim2026kcovdataflow},
the foundation of this work, extends \sancov with two compiler modes,
\texttt{-fsanitize-coverage=trace-args} and \texttt{trace-ret}. The first
reports every source-level parameter of every instrumented frame, inlined
frames included. The second reports every return value. A struct passed by
pointer is expanded field by field through a DWARF-derived
\emph{(offset, size)} table, read with \texttt{copy\_from\_kernel\_nofault()}.
The kernel backend writes variable-length records into a per-task (or
per-remote-handle) buffer exposed at \texttt{/sys/kernel/debug/kcov\_dataflow}.
Each record has a three-word header (sequence, type \texttt{0xE}/\texttt{0xF}/\texttt{0xC},
value count, size, argument index), then the PC, then the source address, then
the values. A field that cannot be read is recorded as the sentinel
\texttt{0xBADADD85}. With \texttt{CONFIG\_KCOV\_ENABLE\_COMPARISONS}, the
\tracecmp operand pairs land in the \emph{same} stream as \texttt{0xC}
records, so argument values and comparison operands of one execution arrive
interleaved and in sequence order. The foundational paper measures
18--21\,ns per stored record and a 66\% syscall-latency cost for whole-kernel
instrumentation with collection idle~\cite{kim2026kcovdataflow}.

\kcovdf provides the signal. This paper is about what a fuzzer has to do to
benefit from it on a real target: deliver valid stateful input, route each
connection's records to the right consumer, fold values into coverage without
flooding the map, turn comparison operands into mutations, and recognize the
bugs that do not crash.

\begin{figure}[t]
\centering
\begin{tikzpicture}[
  node distance=0.3cm,
  box/.style={rectangle, draw, thick, minimum width=2.4cm, minimum height=0.6cm, font=\scriptsize\bfseries, align=center},
  msg/.style={-{Stealth[length=2mm]}, thick},
  lbl/.style={font=\scriptsize, midway, above},
  note/.style={font=\scriptsize\itshape, text=gray},
  rec/.style={-{Stealth[length=1.6mm]}, densely dashed, blue!60!black},
]
\node[box] (client) {Grain harness\\(\libfuzzer, 127.0.0.$n$)};
\node[box, right=4.2cm of client] (server) {\ksmbd\\(conn thread + kworker)};
\node[box, right=1.6cm of server, fill=blue!6] (buf) {kcov\_dataflow\\buffer[$n$]};

% Vertical lifelines
\draw[dashed, gray] (client.south) -- ++(0,-6.0);
\draw[dashed, gray] (server.south) -- ++(0,-6.0);
\draw[dashed, gray] (buf.south) -- ++(0,-6.0);

% Fixed prefix
\begin{scope}[on background layer]
\fill[green!7] ([xshift=-1.4cm, yshift=-0.35cm]client.south) rectangle ([xshift=0.3cm, yshift=-3.85cm]server.south);
\end{scope}
\draw[msg] ([yshift=-0.7cm]client.south) -- node[lbl] {NEGOTIATE} ([yshift=-0.7cm]server.south);
\draw[msg] ([yshift=-1.3cm]client.south) -- node[lbl] {SESSION\_SETUP (NTLMSSP NEGOTIATE)} ([yshift=-1.3cm]server.south);
\draw[msg] ([yshift=-1.9cm]server.south) -- node[lbl] {challenge} ([yshift=-1.9cm]client.south);
\draw[msg] ([yshift=-2.5cm]client.south) -- node[lbl] {SESSION\_SETUP (NTLMv2 AUTH)} ([yshift=-2.5cm]server.south);
\draw[msg] ([yshift=-3.1cm]client.south) -- node[lbl] {TREE\_CONNECT, CREATE} ([yshift=-3.1cm]server.south);
\draw[msg] ([yshift=-3.6cm]server.south) -- node[lbl] {SessionId, TreeId, FileId} ([yshift=-3.6cm]client.south);

% Target PDU
\draw[msg, red!70!black] ([yshift=-4.5cm]client.south) -- node[lbl, text=red!70!black] {target PDU \colorbox{yellow!30}{mutated}} ([yshift=-4.5cm]server.south);
\draw[rec] ([yshift=-4.8cm]server.south) -- node[lbl, text=blue!60!black, font=\tiny] {0xE / 0xF / 0xC records} ([yshift=-4.8cm]buf.south);
\draw[msg] ([yshift=-5.2cm]server.south) -- node[lbl] {NTSTATUS} ([yshift=-5.2cm]client.south);
\draw[rec] ([yshift=-5.6cm]buf.south) -- node[lbl, text=blue!60!black, font=\tiny] {mmap: fold $\to$ ctr[\,], operands $\to$ mutator} ([yshift=-5.6cm]client.south);

% Annotations
\node[note, anchor=east] at ([xshift=-0.15cm, yshift=-2.1cm]client.south) {fixed prefix};
\node[note, anchor=east] at ([xshift=-0.15cm, yshift=-2.45cm]client.south) {(never mutated)};
\draw[decorate, decoration={brace, amplitude=3pt, mirror, raise=2pt}, gray]
  ([xshift=-1.45cm, yshift=-0.45cm]client.south) -- ([xshift=-1.45cm, yshift=-3.75cm]client.south);
\node[note, anchor=east] at ([xshift=-0.15cm, yshift=-4.5cm]client.south) {fuzzed};
\end{tikzpicture}
\caption{A grain. The authenticated prefix (shaded) is a valid, unmutated SMB
exchange; \libfuzzer mutates only the final target PDU. \ksmbd services the
connection on its own threads, and the remote handle derived from the
destination address $127.0.0.n$ routes the resulting \kcovdf records into
buffer~$n$, which only this grain maps.}\label{fig:grain}
\end{figure}

\subsection{Why Existing Approaches Stall}

\paragraph{Syscall-level kernel fuzzing.} \syzkaller and its
descendants~\cite{sun2021healer,wang2021syzvegas,pailoor2018moonshine}
generate sequences of system calls guided by edge coverage. For \ksmbd the
relevant system calls are a TCP \texttt{connect()} and \texttt{sendmsg()} of a
blob, so the SMB state machine becomes an opaque byte string that the fuzzer
has to rediscover. Its coverage signal is value-blind. Its oracle is the
kernel log.

\paragraph{Stateful network-protocol fuzzing.} AFLNet~\cite{pham2020aflnet},
StateAFL~\cite{natella2022stateafl}, SGFuzz~\cite{ba2022sgfuzz} and
Nyx-Net~\cite{schumilo2022nyxnet} model a server's states from message
sequences, response codes, or snapshots. They target user-space servers,
derive state from responses rather than from the server's internal values,
and assume per-process coverage that does not exist for a kernel server whose
work runs on kworkers.

\paragraph{Data-flow-guided fuzzing.} Taint-based fuzzers~\cite{chen2018angora,gan2020greyone}
pay a per-byte shadow cost that is infeasible in the kernel.
IJON~\cite{aschermann2020ijon} requires manual state annotations. RedQueen~\cite{aschermann2019redqueen}
and CmpLog~\cite{fioraldi2020aflpp} obtain input-to-state correspondence from
comparison operands, and Eclipser~\cite{choi2019eclipser} steers values
towards branch conditions. \toolname borrows the input-to-state idea but
draws its operands from the kernel's \tracecmp stream and its argument values
from \kcovdf.

\paragraph{Semantic oracles.} Hydra~\cite{kim2019hydra} checks file systems
for semantic bugs such as crash-consistency violations, and
PeX~\cite{zhang2019pex} and Crix~\cite{lu2019crix} look statically for missing
permission checks. None of them checks a network file server's authorization
contract at run time.

Table~\ref{tab:positioning} summarizes the gap.

\begin{table}[t]
\centering
\caption{Capabilities needed to fuzz \ksmbd's deep, value-dependent,
non-crashing surface.}\label{tab:positioning}
\small
\begin{tabular}{@{}lccccc@{}}
\toprule
& \textbf{Kernel} & \textbf{Valid auth.} & \textbf{Value} & \textbf{Kernel cmp} & \textbf{Integrity} \\
\textbf{Approach} & \textbf{coverage} & \textbf{prefix} & \textbf{feedback} & \textbf{i2s} & \textbf{oracle} \\
\midrule
\syzkaller (syscall level)              & \cmark & \xmark & \xmark & \xmark & \xmark \\
AFLNet / StateAFL / Nyx-Net             & \xmark & \cmark & \xmark & \xmark & \xmark \\
RedQueen / AFL++ CmpLog (user space)    & \xmark & ---    & \xmark & \cmark & \xmark \\
\kcovdf alone (signal only)             & \cmark & \xmark & \cmark & \cmark & \xmark \\
\toolname                               & \cmark & \cmark & \cmark & \cmark & \cmark \\
\bottomrule
\end{tabular}
\end{table}

\subsection{Threat Model and Assumptions}

The attacker is a network client that can reach \ksmbd over TCP port~445 or
over SMBDirect. The attacker holds valid credentials for a \emph{low-privilege}
account, or uses guest access where the configuration allows it, and wants to
(a)~corrupt kernel memory, (b)~deny service, or (c)~violate the integrity of
files the account may not modify. Bugs reachable before authentication are in
scope as well: the \texttt{negotiate} and \texttt{session\_setup} grains fuzz
exactly those PDUs. One reported bug (\texttt{ipc\_validate\_msg()},
Table~\ref{tab:bugs}) requires a malfunctioning or compromised
\texttt{ksmbd.mountd}. That attacker is stronger than the network client, and I
report the bug as kernel hardening against a less-trusted user-space
component. The fuzzer runs entirely inside a \virtme guest on loopback with a
sanitizer-heavy kernel (\kasan, UBSAN, lockdep, \texttt{DEBUG\_LIST}). No
production system or third-party server is involved.


\section{System Design \& Architecture}\label{sec:design}

\subsection{Overview}

\toolname follows one rule throughout: \emph{Python manages, C fuzzes}. The
orchestrator (\texttt{ksmbdzzer.py}) enumerates grains, compiles and launches
harnesses, aggregates their metrics and runs the oracles, and it performs no
per-input work. Each grain is a separate \libfuzzer process. Its
\texttt{LLVMFuzzerTestOneInput} drives one SMB procedure over its own
loopback address, folds that connection's \kcovdf records into
\texttt{\_\_libfuzzer\_extra\_counters}, and returns; its
\texttt{LLVMFuzzerCustomMutator} turns the same records into directed
mutations. \libfuzzer's own corpus logic, which keeps an input that produces
a new feature, is therefore steered entirely by \emph{kernel} behaviour
(Figure~\ref{fig:arch}).

\begin{figure}[t]
\centering
\begin{tikzpicture}[
  font=\scriptsize,
  zone/.style={rectangle, draw, thick, rounded corners=3pt, inner sep=0pt},
  ztitle/.style={font=\scriptsize\bfseries, anchor=north west, inner sep=2pt},
  sub/.style={rectangle, draw, rounded corners=2pt, dashed, inner sep=0pt},
  stitle/.style={font=\tiny\bfseries, anchor=north west, inner sep=1.5pt, text=black!75},
  comp/.style={rectangle, draw, fill=white, minimum height=0.62cm, align=center, font=\tiny, inner sep=1.5pt},
  arrow/.style={-{Stealth[length=1.6mm]}, thick},
  darrow/.style={-{Stealth[length=1.6mm]}, thick, dashed},
  lab/.style={font=\tiny\itshape, fill=white, inner sep=0.8pt, align=center},
]
%% HOST
\node[zone, fill=gray!8, minimum width=16cm, minimum height=1.45cm, anchor=north west] (host) at (0,0) {};
\node[ztitle] at (host.north west) {Host};
\node[comp, minimum width=2.9cm, text width=2.78cm, fill=gray!15] (clang) at (2.1,-0.85) {custom clang/LLVM\\SanCov trace-args/ret/cmp};
\node[comp, minimum width=2.6cm, text width=2.48cm, fill=gray!15] (vng) at (5.25,-0.85) {\texttt{vng --build}\\portable config items};
\node[comp, minimum width=2.4cm, text width=2.28cm, fill=gray!15] (bg) at (8.0,-0.85) {\texttt{build-grains}\\compile cache};
\node[comp, minimum width=2.6cm, text width=2.48cm, fill=gray!15] (cmp) at (10.75,-0.85) {comparison runner\\arms $\times$ trials, cold};
\node[comp, minimum width=2.4cm, text width=2.28cm, fill=gray!15] (db) at (13.5,-0.85) {\texttt{.fuzzdb} (9p)\\durable corpus};
\node[comp, minimum width=0.9cm, text width=0.78cm, fill=gray!15] (con) at (15.4,-0.85) {serial\\console};
\draw[arrow] (clang) -- (vng);
%% GUEST USER SPACE
\node[zone, fill=yellow!6, minimum width=16cm, minimum height=5.0cm, anchor=north west] (us) at (0,-1.7) {};
\node[ztitle] at (us.north west) {Guest user space (\virtme, loopback only)};
% orchestrator
\node[sub, fill=yellow!14, minimum width=8.5cm, minimum height=1.25cm, anchor=north west] (orch) at (0.2,-2.05) {};
\node[stitle] at (orch.north west) {\texttt{ksmbdzzer.py} orchestrator: round loop};
\node[comp, minimum width=1.8cm, text width=1.68cm] (p1) at (1.25,-2.85) {P1 enumerate\\lib + raw + v2};
\node[comp, minimum width=1.8cm, text width=1.68cm] (p2) at (3.3,-2.85) {P2 saturate\\ALIGNED gate};
\node[comp, minimum width=1.8cm, text width=1.68cm] (p3) at (5.35,-2.85) {P3 combine\\shared object};
\node[comp, minimum width=1.8cm, text width=1.68cm] (p4) at (7.4,-2.85) {P4 save\\feed-forward};
\draw[arrow] (p1) -- (p2); \draw[arrow] (p2) -- (p3); \draw[arrow] (p3) -- (p4);
% oracles
\node[sub, fill=red!6, minimum width=7.1cm, minimum height=1.25cm, anchor=north west] (orc) at (8.85,-2.05) {};
\node[stitle] at (orc.north west) {oracles (checks, not phases)};
\node[comp, minimum width=1.9cm, text width=1.78cm] (stab) at (9.95,-2.85) {stability\\$\mathit{ft}_0 < 0.5\cdot$base};
\node[comp, minimum width=2.6cm, text width=2.48cm] (ctr) at (12.3,-2.85) {contract suite: cross-user,\\lease, race, teardown, frag};
\node[comp, minimum width=1.6cm, text width=1.48cm] (arb) at (14.95,-2.85) {arbiter\\\texttt{fb.bin}};
% scheduler
\node[comp, minimum width=15.6cm, text width=15.48cm, fill=orange!8] (sched) at (8.0,-3.75) {\texttt{gen.py run\_grains()}: slot $s \mapsto 127.0.0.(2{+}s \bmod 250)$ $\cdot$ replicas share a corpus $\cdot$ saturation detection $\cdot$ fast-bail $\cdot$ PRODUCTIVE gate $\cdot$ fleet-union PC bitmap};
% grain fleet
\node[sub, fill=orange!12, minimum width=13.9cm, minimum height=2.15cm, anchor=north west] (fleet) at (0.2,-4.3) {};
\node[stitle] at (fleet.north west) {grain fleet: one \libfuzzer process per grain or replica, each on its own 127.0.0.$n$};
\node[comp, minimum width=2.5cm, text width=2.38cm] (lib) at (1.6,-5.05) {library tier (most grains)\\\texttt{GRAINS[]} + libsmbclient};
\node[comp, minimum width=2.3cm, text width=2.18cm] (raw) at (4.15,-5.05) {raw tier\\\texttt{common.h} + NTLMv2};
\node[comp, minimum width=1.7cm, text width=1.58cm] (v2) at (6.25,-5.05) {v2 hybrid\\i2s at stuck};
\node[comp, minimum width=1.9cm, text width=1.78cm] (cmb) at (8.2,-5.05) {combo harness\\$A{+}B$ shared fid};
\node[comp, minimum width=2.3cm, text width=2.18cm, fill=cyan!10] (fold) at (10.6,-5.05) {fold $\to$ \texttt{ctr[4096]}\\path + first-seen value};
\node[comp, minimum width=2.0cm, text width=1.88cm, fill=red!8] (mut) at (12.95,-5.05) {\texttt{mutate\_i2s}\\cmp ring, ret tokens};
\draw[arrow] (fold) -- (mut);
\node[comp, minimum width=1.6cm, text width=1.48cm, fill=gray!12] (mountd) at (15.05,-5.05) {\texttt{ksmbd.mountd}\\users, shares};
%% GUEST KERNEL
\node[zone, fill=green!6, minimum width=16cm, minimum height=3.05cm, anchor=north west] (kern) at (0,-6.95) {};
\node[ztitle, anchor=south west] at (kern.south west) {Guest kernel (linux-next, whole kernel built with trace-args/ret/cmp)};
\node[comp, minimum width=2.6cm, text width=2.48cm, fill=green!12] (tcp) at (1.6,-7.75) {TCP :445, handle\\\texttt{0x4b44}$\,|\,$low16(daddr)};
\node[comp, minimum width=2.3cm, text width=2.18cm, fill=green!12] (rdma) at (4.25,-7.75) {SMBDirect (RXE/SIW)\\5 annotated work fns};
\node[comp, minimum width=2.7cm, text width=2.58cm, fill=green!12] (conn) at (7.0,-7.75) {conn thread\\\texttt{ksmbd\_conn\_handler\_loop}};
\node[comp, minimum width=3.7cm, text width=3.58cm, fill=green!12] (work) at (10.45,-7.75) {kworker \texttt{handle\_ksmbd\_work}\\$\to$ 19 SMB2 handlers $\to$ VFS};
\node[comp, minimum width=1.8cm, text width=1.68cm, fill=green!12] (ipc) at (14.85,-7.75) {\texttt{transport\_ipc}\\generic netlink};
\node[comp, minimum width=3.6cm, text width=3.48cm, fill=blue!8] (rem) at (6.0,-9.0) {remote sections\\\texttt{kcov\_df\_remote\_start\_common}};
\node[comp, minimum width=3.8cm, text width=3.68cm, fill=blue!8] (kdf) at (10.45,-9.0) {\kcovdf: remote map, per-worker scratch, atomic merge $\to$ buffer[$n$]};
\node[comp, minimum width=2.4cm, text width=2.28cm, fill=red!8] (san) at (14.1,-9.0) {\kasan, UBSAN, lockdep,\\\texttt{DEBUG\_LIST}};
\draw[arrow] (rdma) -- (conn); \draw[arrow] (tcp.north) -- ++(0,0.22) -| ([xshift=-0.4cm]conn.north); \draw[arrow] (conn) -- (work);
\draw[arrow] (conn.south) -- (conn.south |- rem.north); \draw[arrow] (work.south) -- node[lab] {remote section} (work.south |- kdf.north);
\draw[arrow] (rem) -- (kdf);
%% CROSS-ZONE FLOWS
\draw[arrow] (p2.south) -- node[lab, right=1pt] {launch} (p2.south |- sched.north);
\draw[arrow] (lib.south) -- node[lab, pos=0.45] {SMB2 PDUs to\\127.0.0.$n$:445} (lib.south |- tcp.north);
\draw[darrow] (raw.south) -- node[lab, pos=0.4] {\texttt{rdma}\\grain} (rdma.north);
\draw[arrow, blue!60!black] (kdf.east) -- (12.55,-9.0) -- (12.55,-6.2) -| node[lab, pos=0.25, text=blue!60!black] {mmap: this connection's records} (fold.south);
\draw[{Stealth[length=1.6mm]}-{Stealth[length=1.6mm]}, thick] (mountd.south) -- node[lab] {login,\\tree} (mountd.south |- ipc.north);
\draw[arrow, red!60!black] (san.east) -- (15.92,-9.0) |- node[lab, pos=0.15, text=red!60!black, rotate=90] {splats} (con.east);
\draw[arrow] (p4.north) -- (7.4,-1.58) -| node[lab, pos=0.25] {corpus, features, value pool} (db.south);
\end{tikzpicture}
\caption{System architecture of \toolname. On the host, a custom clang builds
the kernel with \traceargs/\traceret/\tracecmp, and harnesses are pre-compiled
into a cache. In the guest, the orchestrator runs the four-phase round loop and
the oracles, and \texttt{gen.py} schedules a fleet of \libfuzzer grain
processes, each dialling its own loopback address. In the kernel, \ksmbd's
transports derive a remote handle from that address, and the connection thread,
kworkers and SMBDirect work functions run in remote sections, so \kcovdf
merges the argument, return and comparison records into the buffer that only
that grain maps. Sanitizer reports leave the guest on the serial console, and
corpora are mirrored to a host-durable store after every round.}\label{fig:arch}
\end{figure}

\subsection{Grains: Deep by Construction}\label{sec:grain}

A \emph{grain} is a working, ordinary SMB scenario, written down as code. It
performs the prerequisite exchange (negotiate, authenticate, tree connect,
and usually \texttt{CREATE}) with valid values, and then issues one
\emph{target} procedure whose fields or parameters come from the fuzzer
input. The prefix is never mutated. This is the central structural decision:
mutation begins \emph{past} the parser and the session checks, so no budget
is spent rediscovering the handshake, and every input exercises handler code
on its first execution.

\paragraph{Two tiers.} Most grains (the \emph{library} tier) are functions in
the \texttt{GRAINS[]} registry of \texttt{libksmbdzzer.so}. They obtain the
prefix from \texttt{libsmbclient}, or from a raw, already-authenticated
connection pool, and use the input to parameterize the target operation:
offsets, lengths, information classes, FSCTL buffers, create options,
compound layouts. A smaller \emph{raw} tier of socket harnesses speaks SMB2
directly: it builds the 64-byte sync header and the 4-byte NetBIOS length
prefix itself, authenticates with its own NTLMv2 implementation, and lets the
input \emph{be} the target PDU. The library tier inherits correct signing and
session handling. The raw tier can corrupt structure sizes, offsets and
compound links that \texttt{libsmbclient} would never emit. Each tier covers
the other's blind spot.

\paragraph{Coverage of the procedure surface.} The registry defines 211
grains and runs 210 of them; one ACL grain is suppressed because another
grain reaches the same \texttt{parse\_sec\_desc()} surface. The grains are
ordered by kernel data-path depth, with file and VFS I/O first and session
and authentication last, so that the grain index means the same thing in the
documentation, the registry and the compile order. Table~\ref{tab:surface}
shows how the fleet maps onto the protocol. Beyond one grain per handler,
112 grains target parser depth and interaction: chained compound requests,
array walkers left at one iteration by single-element grains, durable and
lease reconnect races, credit exhaustion, transform nesting, and partial or
oversize framing.

\begin{table}[t]
\centering
\caption{Procedure surface covered by the 210 active grains (canonical
grain-to-handler map in \texttt{GRAIN.md}).}\label{tab:surface}
\small
\begin{tabular}{@{}lll@{}}
\toprule
\textbf{Surface} & \textbf{Coverage} & \textbf{Example grains} \\
\midrule
SMB2 commands (\texttt{smb2pdu.c})     & 19 of 19, each with a dedicated grain & \texttt{write}, \texttt{lock}, \texttt{notify}, \texttt{oplock\_ack} \\
FSCTL (\texttt{smb2\_ioctl})           & all 16 handled codes                 & \texttt{copychunk}, \texttt{fsctl\_zero}, \texttt{fsctl\_dupext} \\
\texttt{SET\_INFO}/\texttt{QUERY\_INFO} & file, FS and security classes       & \texttt{rename}, \texttt{hardlink}, \texttt{set\_alloc} \\
\texttt{CREATE} contexts               & 11 contexts                          & \texttt{create\_dh2}, \texttt{create\_aapl}, \texttt{lease\_v2} \\
SMB3 features                          & encryption, signing, multichannel, durable-v2 & \texttt{encrypt}, \texttt{session\_bind}, \texttt{dh2c} \\
Transport                              & TCP framing, SMBDirect (RDMA)        & \texttt{partial\_pdu\_dribble}, \texttt{rdma} \\
SMB1 compatibility                     & generic + 5 per-opcode grains        & \texttt{smb1\_ntcreate}, \texttt{smb1\_trans} \\
IPC / RPC                              & named pipes, DCE/RPC opnums, netlink & \texttt{ndr}, \texttt{rpc\_opnum}, \texttt{ipc} \\
Parser depth / interaction             & 112 grains                           & \texttt{compound\_related\_fid}, \texttt{credit\_exhaust} \\
\bottomrule
\end{tabular}
\end{table}

\paragraph{Five archetypes.} Although there are 211 grains, they follow only
five procedure shapes (Table~\ref{tab:archetypes}). Every grain shares one
skeleton: it arms the per-connection coverage once, establishes its preamble,
and for each input resets the buffer, sends its fuzzed PDU, and folds the
resulting records. Grains differ in two respects only: \emph{what the fuzzed
endpoint is}, and \emph{how many messages it takes to reach it}. The archetype
also determines how a grain recovers when an input kills its connection,
which is what keeps a long campaign productive.

\begin{table}[t]
\centering
\caption{Grain archetypes, distilled from \texttt{GRAIN.md}. Counts are
approximate shares of the 211-grain registry.}\label{tab:archetypes}
\footnotesize
\setlength{\tabcolsep}{4pt}
\newcolumntype{L}[1]{>{\raggedright\arraybackslash}p{#1}}
\begin{tabular}{@{}L{2.1cm}L{4.4cm}L{4.1cm}L{0.9cm}L{3.9cm}@{}}
\toprule
\textbf{Archetype} & \textbf{Shape (prefix $\to$ fuzzed endpoint)} & \textbf{What it stresses} & \textbf{Share} & \textbf{Representative grains} \\
\midrule
\textbf{P} single-op & authenticated preamble and a persistent \texttt{FileId}; each input is \emph{one} operation on it & handler argument decoding, bounds and offset checks, VFS calls & $\approx$150 & \texttt{write\_ext}, \texttt{fsctl\_*}, \texttt{set\_*}, \texttt{create\_*} \\
\textbf{M} multi-op & several ordered operations per input: compound chain, deferred lock with cancel, close during an operation & cross-operation state, teardown ordering, race windows & $\approx$40 & \texttt{compound\_related\_fid}, \texttt{write\_lock\_race}, \texttt{durable\_reconnect\_race} \\
\textbf{R} RPC / pipe & \texttt{IPC\$} tree, pipe open, DCE/RPC \texttt{BIND}; each input is one RPC request & NDR decoding, opnum dispatch, pipe transceive & $\approx$7 & \texttt{ndr}, \texttt{rpc\_opnum}, \texttt{pipe\_transceive\_bind} \\
\textbf{D} disrupting & negotiate only; the fuzzed PDU is session- or protocol-level and may tear the connection down, so the grain reconnects & SPNEGO/NTLMSSP, signing, transform headers, SMB1, logoff, framing & $\approx$30 & \texttt{spnego\_asn1}, \texttt{encrypt}, \texttt{logoff\_*}, \texttt{smb1\_*} \\
\textbf{X} architectural & reaches \ksmbd through a path without a per-connection IP handle & netlink IPC to \texttt{ksmbd.mountd}, SMBDirect data plane & 2 & \texttt{ipc}, \texttt{rdma} \\
\bottomrule
\end{tabular}
\end{table}

\paragraph{Alignment.} A grain is only as good as its prefix. A wrong
\texttt{StructureSize}, a dialect rejected early, or a pool connection torn
down by an earlier grain makes \ksmbd drop the PDU before any handler runs,
yet the grain keeps running and keeps consuming budget. \toolname enforces
alignment twice. Before a campaign, \texttt{selftest} runs each grain several
times and classifies it by the number of distinct kernel PCs its connection
produced: WORKS ($\geq$1 PC), DEAD (it ran but no PC), BAIL (a prerequisite
failed), or EXCL (architecturally unmeasurable on a per-connection handle).
The measurement uses kernel coverage, not the SMB reply, because an error
reply is produced both by a PDU rejected at dispatch and by one that ran the
whole handler and then failed a check. During a campaign, \texttt{gen.py}
tags a grain ALIGNED only if its \libfuzzer feature count reaches
$\mathit{ft} \geq 100$. In practice deep grains reach thousands of features
and misaligned ones stay below~30.

\subsection{Per-Connection Kernel Coverage Routing}\label{sec:route}

The fleet runs many grains at once, and \ksmbd services all of them on shared
kworkers. A buffer that mixed every grain's records would credit one grain's
input with another grain's coverage. \toolname gives each worker its own
loopback destination, $127.0.0.n$, and routes by destination:
\begin{itemize}[nosep]
  \item In \texttt{ksmbd\_conn\_alloc()} every connection receives a default
        \texttt{kcov\_common\_handle()}. For an IPv4 TCP connection,
        \texttt{alloc\_transport()} overrides it with
        \[ \mathit{handle}(c) = \texttt{0x4b440000} \;\big|\; \big(\mathrm{ntohl}(\texttt{inet\_rcv\_saddr}(c)) \wedge \texttt{0xFFFF}\big), \]
        that is, the \emph{local} (destination) address the client dialled.
        The \texttt{0x4b44} (``KD'') tag keeps the value out of the
        small-integer handle space that other subsystems use.
  \item The connection handler loop and the request kworker
        (\texttt{handle\_ksmbd\_work()}) bracket their bodies with
        \texttt{kcov\_remote\_start\_common()}/\texttt{kcov\_df\_remote\_start\_common()}
        on that handle, and with the matching stops. Five SMBDirect work
        functions (negotiate receive on accept and connect, the idle timer,
        immediate send, and receive-buffer refill) are wrapped the same way
        through thin wrappers around their bodies, so that early returns
        cannot unbalance a section. The RDMA transport copies the
        connection's handle into its SMBDirect socket.
  \item The grain process computes the same value from its own octet, maps
        a \kcovdf buffer, and publishes the handle with
        \texttt{KCOV\_DF\_REMOTE\_ENABLE}. Remote writers assemble records in
        a per-worker scratch area and merge them into the shared buffer with
        one atomic cursor update~\cite{kim2026kcovdataflow}.
\end{itemize}
The design costs a few lines in \ksmbd's transport and nothing on the hot
path. It requires no change to the SMB protocol and no cooperation from the
client library, since \texttt{libsmbclient} only needs to be told which
address to dial. The low 16 bits of the address name the worker, so up to
250 concurrent workers fit in $127.0.0.\{2..251\}$. IPv6 connections remain
unrouted.

\subsection{Value-Aware Coverage Folding}\label{sec:fold}

Each execution leaves a sequence of records
$r = \langle \tau, \mathit{pc}, a, v_0..v_{m-1}\rangle$ in the buffer, where
$\tau \in \{\texttt{E}, \texttt{F}, \texttt{C}\}$ and $m$ is the number of
values. The harness must turn these records into \libfuzzer features. Folding
every $(\mathit{pc}, v)$ pair naively makes the map unbounded: every kernel
pointer, sequence number or timestamp becomes a ``new'' feature, the corpus
grows without limit, and a grain never saturates, which removes the very
signal the scheduler uses (Section~\ref{sec:loop}). Folding only the PC
discards what \kcovdf was built to provide. \toolname splits the 4096-entry
counter table into two regions (Figure~\ref{fig:fold}). The description below
is the raw-tier implementation. The library tier implements the same policy
by feature hashing (Section~\ref{sec:impl}).

\paragraph{Path region.} The path region,
$\mathit{ctr}[0..P{-}1]$ with $P = 3584$ by default, is folded from the PC
alone, $\phi_\mathrm{path}(r) = \lfloor (\mathit{pc}\cdot c_1)/2^{12}\rfloor \bmod P$.
It is bounded by the set of reachable PCs, so it \emph{saturates}. Saturation
is a feature here, not a defect: it is the clean ``no new behaviour'' signal
on which grain scheduling depends.

\paragraph{Bounded value region.} For every entry or return record and every
field $k$ whose value is neither zero nor the unreadable sentinel, the harness
computes a 16-bit bucket
\[ b_k = \Big\lfloor \big((\mathit{pc} \oplus \kappa(v_k)) + k\cdot c_2\big)\cdot c_3 \,/\, 2^{48} \Big\rfloor \]
and increments $\mathit{ctr}[P + (b_k \bmod V)]$ \emph{only the first time}
$b_k$ is seen, tracked in a persistent $2^{16}$-bit bitmap. Here $V = 512$,
$c_3$ is the 64-bit golden-ratio constant, and mixing in the field index $k$
keeps the same value in different struct fields distinct. A loop that reads
the same length a million times produces one feature, not a million. A
request that drives a length, flag word or create option to a value
\emph{never before seen at that call site} is rewarded and kept.

\paragraph{Value classes.} Raw values still over-reward noise: two allocations
give two pointers, and lengths of 1000 and 1001 behave identically. The
canonical arm therefore folds a value class
\[ \kappa(v) = \begin{cases}
0 & v = 0\\
\texttt{0x1000} \;|\; (v \wedge \texttt{0x38}) & v \geq \texttt{0xffff}\cdot 2^{48} \quad\text{(kernel pointer)}\\
\texttt{0x2000} & v > 2^{47}-1 \quad\text{(other non-canonical)}\\
\lfloor \log_2 v \rfloor & \text{otherwise (scalar magnitude)}
\end{cases} \]
which keeps NULL-ness, a little alignment, and order of magnitude, so that
crossing a boundary such as $2^{31}$ or $2^{32}$ is still a new feature. A
relational variant adds, for each pair of fields $(a, b)$ in one record, the
three-way ordering $\rho(v_a, v_b) \in \{<, =, >\}$. Value classes collapse
an offset and a length of the same magnitude into one bucket, and the ordering
splits them again, because whether \texttt{offset < length} holds is exactly
what decides an out-of-bounds access.

\paragraph{Comparisons are not coverage.} Comparison records reach the path
region only through their PC. Their operands are usually the fuzzer's own
input bytes echoed back, so folding them would mint a new feature on almost
every mutation, flood the value region, and drown the genuine argument
values. Of the two ways to use comparisons, \emph{operand as coverage}
(libFuzzer's value profile) or \emph{context plus directed solving},
\toolname takes the second: the PC refines the location, and the operand
pair goes to the mutator (Section~\ref{sec:i2s}).

\begin{figure}[t]
\centering
\begin{tikzpicture}[
  cell/.style={rectangle, draw, minimum height=0.7cm, font=\scriptsize, align=center},
  box/.style={rectangle, draw, rounded corners, fill=gray!8, minimum height=0.6cm, font=\scriptsize, align=center},
  note/.style={font=\scriptsize\itshape, text=black!75, fill=white, inner sep=1pt},
  arrow/.style={-{Stealth[length=1.6mm]}, thick},
]
\node[cell, fill=orange!20, minimum width=10.5cm] (path) at (0,0) {PATH region: \texttt{ctr[0..3583]}, folded from $\phi(\mathit{pc})$, saturates};
\node[cell, fill=cyan!20, minimum width=3.2cm, anchor=west] (val) at (path.east) {VALUE region: \texttt{ctr[3584..4095]}};
\node[box] (rec) at (0,1.9) {record $\langle \tau, \mathit{pc}, a, v_0..v_{m-1} \rangle$, $\tau \in \{\texttt{0xE}, \texttt{0xF}, \texttt{0xC}\}$};
\node[box, fill=red!8] (mut) at ([xshift=-3.6cm]rec.west |- 0,1.9) {\texttt{mutate\_i2s}\\operand ring};
\draw[arrow] (rec.south -| path.north) -- node[note] {every record: $\mathit{pc}$} (path.north);
\draw[arrow] (rec.south east) -- node[note, pos=0.45] {\texttt{0xE}/\texttt{0xF} only: first-seen $(\mathit{pc}, k, \kappa(v_k))$} (val.north);
\draw[arrow, dashed] (rec.west) -- node[note, above=2pt] {\texttt{0xC} operands} (mut.east);
\end{tikzpicture}
\caption{The coverage fold. The path region gives edge coverage's bounded,
saturating signal. The value region adds one feature per never-before-seen
(call site, field, value class) triple. Comparison operands drive mutation
instead of coverage.}\label{fig:fold}
\end{figure}

\paragraph{Engine arms.} Because the fold policy lives in the harness and not
in the kernel, the same grains on the same kernel can be run as several
engines, selected by \texttt{KSMBDZZER\_ENGINE}, with identical
instrumentation (Table~\ref{tab:arms}). Comparing \texttt{dataflow-vec} with
\texttt{pc-i2s} isolates the payoff of value coverage at the same
per-execution cost, and comparing \texttt{pc-i2s} with \texttt{pc-havoc}
isolates the payoff of input-to-state mutation. A sixth configuration replaces
\kcovdf with mainline KCOV (\texttt{KSMBDZZER\_KCOV=1}) on the same
per-connection handle, as a stock edge-coverage baseline.

\begin{table}[t]
\centering
\caption{Engine arms. All arms use the same kernel and the same grains; only
the fold and the mutator differ.}\label{tab:arms}
\small
\begin{tabular}{@{}lll@{}}
\toprule
\textbf{Arm} & \textbf{Coverage folded from the record stream} & \textbf{Mutator} \\
\midrule
\texttt{dataflow}     & path + raw value of the record              & i2s + havoc \\
\texttt{dataflow-vec} & path + every field, value-class normalized (canonical) & i2s + havoc \\
\texttt{dataflow-rel} & \texttt{dataflow-vec} + pairwise field ordering & i2s + havoc \\
\texttt{pc-i2s}       & path only (values dropped)                  & i2s + havoc \\
\texttt{pc-havoc}     & path only (values dropped)                  & havoc \\
\midrule
\texttt{kcov}         & mainline \tracepc, path only                & havoc (no records for i2s) \\
\bottomrule
\end{tabular}
\end{table}

\subsection{Dataflow-Directed Mutation}\label{sec:i2s}

The custom mutator turns the same record stream into mutations
(Algorithm~\ref{alg:i2s}, Appendix~\ref{app:i2s}). Two of every three calls
are plain \libfuzzer havoc, which keeps exploration cheap; the third is
directed. In the hybrid (v2) grains, the directed call is made only at a
saturation point, after 64 consecutive executions without a new path. The
directed step has three passes.

\paragraph{Pass 0: comparison operands (RedQueen).} For each \texttt{0xC}
record of width $w \geq 2$ with distinct operands $(a_1, a_2)$, if the
little-endian encoding of one operand occurs in the input, it is overwritten
with the other, which flips a magic-value check in a single step. In the
library tier this in-place case is rare. There the input only
\emph{parameterizes} a \texttt{libsmbclient} operation, and the operands live
in the PDU that the library builds, not in the bytes \libfuzzer holds. To
handle this, every operand seen by the fold, except trivial values
($\leq 1$) and kernel pointers, is recorded in a persistent, deduplicated
512-entry ring. Up to six operands, drawn alternately from the current
execution and from the ring, are \emph{injected} at seed-chosen offsets of an
input that is first grown to at least 16 bytes. The magic value thereby enters
the bytes that parameterize the PDU, and it can reach the comparison on the
next execution. Without the ring, the operand supply would depend on what the
buffer happens to contain when the mutator runs, which is frequently an
execution with no comparison record at all.

\paragraph{Pass 1--2: arguments to boundaries, returns to arguments.} For each
\texttt{0xE} field value $v$ that occurs in the input (widest match first, at
8, 4 and 2 bytes, since lengths and offsets live in the wide fields), the
mutator replaces it with one of three things: a kernel \emph{return} value
taken from a persistent 512-entry return dictionary or from the current
execution, a value carried over from the previous round, or a type boundary
($0$, $1$, $2^{31}{-}1$, $2^{31}$, $2^{32}{-}1$, $2^{63}{-}1$, $2^{63}$,
$2^{64}{-}1$, $\texttt{0x1000}$, and narrower equivalents). Using returns as
tokens is the \emph{reverse} direction of input-to-state correspondence: a
handle, size or status that the kernel computed at iteration $N$ is exactly
the value a later request will be checked against.

\paragraph{Cross-round arbiter.} After each round, the orchestrator distils
the body offsets that increased coverage or tripped an oracle, together with
the interesting values, into a packed blob (\texttt{fb.bin}, magic
\texttt{0xF00DDA7A}, 64 offsets and 128 values). The raw grains load it at
start-up and include those offsets and values in the directed pass. This
closes a data-driven loop with no recompilation: the harness is compiled
once, and its behaviour comes from data.

The mutator uses no SMT solver, no symbolic execution and no byte-level
taint, none of which is feasible across the user/kernel boundary at fuzzing
speed. It relies on the observation that \ksmbd copies most wire fields into
kernel arguments with little or no transformation, so a value that appears in
the record stream can usually be located in the PDU.

\subsection{Procedure Combination on a Shared Object}\label{sec:combine}

Single-procedure grains cannot find bugs that need two procedures to
interact: a lease taken by one request and abused by another, a byte-range
lock left behind, a handle torn down halfway. Phase~P3 therefore runs grains
in pairs, \texttt{pfz\_grain\_combo2(a, b)}, on \emph{one shared object}. A
file \texttt{combo\_obj} is opened on the pool connection(s) so that both
grains act on the same handle and path, and the first input byte selects one
of three modes:
\begin{enumerate}[nosep]
  \item \textbf{Shared}: $A$ then $B$ on the same connection and handle, with
        no reset in between, so $B$ inherits $A$'s state.
  \item \textbf{Interleave}: $A, B, A, B$ on the shared handle.
  \item \textbf{Race}: $A$ on connection~0 and $B$ on connection~1, at the
        same time. $B$ runs in a forked child that rebinds its connection
        references to connection~1's own socket, so the bytes on the wire
        never interleave, and a five-second alarm bounds a hung read.
\end{enumerate}
All pairs are generated from one compiled combination harness plus one
symbolic link per pair, so the $O(n^2)$ space costs one compilation.

\subsection{Oracles as Checks, not Phases}\label{sec:oracles}

\paragraph{Crash.} \kasan, UBSAN, lockdep, \texttt{DEBUG\_LIST},
refcount-saturation and fault messages on the serial console, matched against
a fixed pattern list. \texttt{panic\_on\_oops} and \texttt{hung\_task\_panic}
turn a wedge into a VM exit, so the run harness can resume the campaign from
the durable corpus.

\paragraph{Stability.} Every round re-runs each grain's bare scenario, and the
first $\mathit{ft}$ that \libfuzzer reports, $\mathit{ft}_0$, is the coverage
of that unmutated scenario. The orchestrator keeps a per-grain baseline that
only ratchets upwards. If a grain with baseline $\geq 100$ now reaches less
than half of it, the \emph{normal} scenario has stopped working. The previous
round's fuzzing has then left \ksmbd in a corrupted but still running state:
a non-crashing bug whose reproducer is in that grain's corpus history.

\paragraph{Contract (write-side integrity).} Once per round, the
\texttt{DataflowDirector} authenticates to a dedicated \texttt{[privtest]}
share, where the fuzzing account maps to a non-root UID, and runs a battery of
probes. Each probe rebuilds its own clean session.
\begin{itemize}[nosep]
  \item \emph{Cross-user integrity.} A victim file owned by \texttt{root:root}
        with mode \texttt{0644}, in a \texttt{0755} root-owned directory,
        holds a known secret. Three procedures are attempted through a
        handle opened \emph{read-only}: \texttt{CREATE} with
        \texttt{OVERWRITE\_IF}, \texttt{CREATE} with
        \texttt{FILE\_DELETE\_ON\_CLOSE}, and a direct \texttt{SMB2\_WRITE}.
        The verdict is taken \emph{from the backing file system}: did the
        secret change, or did the file disappear? It is not taken from the
        returned status, because \ksmbd can return \texttt{STATUS\_SUCCESS}
        while POSIX silently prevents the effect. A hit therefore means the
        server acted with more authority than the user has. A sanity probe
        first confirms that the account is in fact non-root on the share;
        otherwise the oracle declares itself inconclusive.
  \item \emph{Lease grant.} The granted lease state must be a subset of the
        requested state and internally consistent (write caching implies
        read caching), as required by MS-SMB2~\S3.3.5.9.11~\cite{mssmb2}.
  \item \emph{Race integrity.} A known pattern is written, an oplock race is
        driven across pool connections, and the file is read back. A silent
        mismatch is state corruption without a crash.
  \item \emph{Teardown lifetime.} A handle is opened, its owner is destroyed
        (\texttt{TREE\_DISCONNECT} frees the tree connection; \texttt{LOGOFF}
        frees the session with all its trees and files), and the dangling
        \texttt{FileId} is used. \texttt{STATUS\_SUCCESS} is a logic bug; a
        \kasan splat between the printed begin and end markers is a memory
        bug.
  \item \emph{Fragmented framing.} A \texttt{CREATE}/\texttt{WRITE}/\texttt{CLOSE}
        cycle is re-sent with each PDU split across eight TCP segments, so
        that the partial-read loop in \texttt{ksmbd\_tcp\_readv()}, which a
        single loopback write coalesces away, actually runs.
\end{itemize}
Every round also runs a positive control, so that silence has a meaning. A
predicate self-test checks that a known read-only plus \texttt{OVERWRITE}
tuple is flagged and a read-write tuple is not. A live control then opens the
victim with \texttt{FILE\_READ\_DATA} only and writes through that handle;
\ksmbd must reject the write on its granted-access check. Without the
control, a silent oracle could mean either a correct server or a broken
probe.

\subsection{The Round Loop and Feed-Forward}\label{sec:loop}

A campaign is a sequence of \emph{rounds} (generations), each with four
phases.
\begin{description}[nosep, font=\normalfont\bfseries]
  \item[P1 enumerate.] Build the grain set (library, raw and hybrid
        harnesses). A compile cache keyed on the source hash and dependency
        modification times skips \texttt{clang} for an unchanged harness.
        The value pool travels in \libfuzzer's \texttt{-dict} file, not in
        the C source, so carrying values forward never forces a recompile.
  \item[P2 saturate.] Run every grain until its coverage saturates or its time
        cap expires. Saturation means that the relative growth of
        $\mathit{ft}$ per five-second window stays below a ratio (default
        $0.02$) for three consecutive windows, after a warm-up of at least 32
        executions so that the directed mutator has had time to act.
  \item[P3 combine.] All-pairs combination on a shared object
        (Section~\ref{sec:combine}), in the final round, under an explicit
        combination budget.
  \item[P4 save.] Persist each grain's corpus, its features and the value
        pool, and mirror them to a host-durable directory, so that a VM
        wedge costs one round rather than the campaign. The next round
        replays this corpus, so its $\mathit{ft}_0$ starts higher.
\end{description}

Two selection rules keep the fleet honest. \emph{Anti-monkey selection}: a
grain counts as PRODUCTIVE only if mutation beats the bare scenario,
$\Delta = \mathit{ft} - \mathit{ft}_0 \geq 100$ or
$\Delta \geq 0.2\,\mathit{ft}_0$. A grain that reaches deep code but gains
nothing from fuzzing is a monkey at a keyboard and is not ranked or combined.
\emph{Fast-bail}: a grain whose session never came up runs the harness
entirely in-process, at execution rates that no grain making SMB round trips
can reach. A grain that exceeds 20{,}000 executions per second after 50{,}000
executions while staying below the alignment threshold is killed at once,
instead of spending its whole budget producing no kernel coverage.

\section{Implementation}\label{sec:impl}

\subsection{Components}

Table~\ref{tab:loc} lists the components. The kernel side consists of
\kcovdf~\cite{kim2026kcovdataflow} (\texttt{kernel/kcov\_dataflow.c}, 1{,}174
lines) and the remote-coverage annotations in \texttt{fs/smb/server} and
\texttt{fs/smb/smbdirect}, carried as four patches: mainline KCOV and \kcovdf
sections for the server and for SMBDirect.

\begin{table}[t]
\centering
\caption{Implementation size (lines, \texttt{wc -l}).}\label{tab:loc}
\small
\begin{tabular}{@{}lrl@{}}
\toprule
\textbf{Component} & \textbf{LOC} & \textbf{Role} \\
\midrule
\texttt{libksmbdzzer.c} / \texttt{.h} & 6{,}429 / 198 & \texttt{GRAINS[]} registry, \kcovdf mmap + fold, i2s, combo, RDMA \\
\texttt{grain/gen.py}                 & 2{,}886       & harness generator, compile cache, fleet scheduler, metrics \\
\texttt{grain/common.h}               & 1{,}005       & raw-tier harness: SMB2 framing, fold, \texttt{mutate\_i2s} \\
\texttt{grain/ntlmv2.h}               & 92            & raw-tier NTLMv2 authentication \\
\texttt{ksmbdzzer.py}                 & 2{,}483       & orchestrator: round loop, oracles, arbiter, \texttt{init}/\texttt{selftest} \\
\texttt{engine\_compare\_campagin.sh} & 628           & build + multi-arm, multi-trial comparison runner \\
\midrule
\texttt{kernel/kcov\_dataflow.c}      & 1{,}174       & \kcovdf backend (foundational work) \\
\bottomrule
\end{tabular}
\end{table}

\subsection{Two Fold Back-Ends, One Policy}

The raw tier implements the fold of Section~\ref{sec:fold} directly over the
4096-byte \texttt{\_\_libfuzzer\_extra\_counters} section: path region,
first-seen value region, and an optional relational pass. The library tier,
which hosts most of the fleet, hashes each kernel record with 64-bit FNV-1a
into one 32-bit feature. The hash covers the PC alone for comparison records
and for the PC-only arms; the PC and every value-classed field, tagged with
its index, for \texttt{dataflow-vec}; and the PC and the classed value for
returns. The generated harness increments
$\mathit{ctr}[\mathit{feature} \wedge 4095]$. Only records whose PC lies in
the kernel text range ($\geq$\texttt{0xffffffff80000000}) are counted. Both
back-ends implement the same policy: comparisons never enter the value space,
pointers are classed, and the PC-only arms drop values completely, so a
comparison between arms is a controlled ablation in both tiers. A grain's
whole \texttt{LLVMFuzzerTestOneInput} is three calls: run grain $i$, fetch
its features, write the counters (Listing~\ref{lst:grain},
Appendix~\ref{app:harness}).

\subsection{Fleet Scheduling}

\texttt{run\_grains()} assigns slot $s$ the address
$127.0.0.(2 + s \bmod 250)$ and runs grains in waves of \texttt{-procs}
(default: all CPUs) with \texttt{-max\_len=512}, the per-grain dictionary,
and a seed derived from \texttt{KSMBDZZER\_SEED} and the slot. When there
are fewer grains than cores, the idle cores run \emph{replicas} of the same
grain that share its corpus directory. Several engineering measures emerged
from running unattended multi-hour campaigns against a sanitizer kernel:
\begin{itemize}[nosep]
  \item \textbf{Lazy pool authentication.} Only grains that need the
        authenticated two-connection pool create it. When every grain
        authenticated up front, the single \texttt{ksmbd.mountd} IPC daemon
        was flooded with NTLMv2 logins. In the self-test this storm alone
        reduced WORKS from 208 to 78.
  \item \textbf{Session drain.} A wave with many grains that executed zero
        inputs means \ksmbd hit its connection or session ceiling, so the
        scheduler sleeps $\min(8\,\mathrm{s}, 1\,\mathrm{s} \times \text{deaths})$
        before the next wave. \ksmbd is restarted only when a liveness probe
        on port~445 fails, never as a matter of routine.
  \item \textbf{Staggered teardown.} Flat grains saturate on a purely
        time-based schedule, so a whole wave reaches \texttt{terminate()}, and
        closes its \kcovdf descriptors, in the same instant. A per-grain
        delay spreads these closes out. Each grain is a session leader and is
        killed by process group, so forked race children cannot survive while
        holding a \ksmbd socket.
  \item \textbf{Metrics that survive \texttt{SIGKILL}.} Saturation ends a grain
        with a signal that bypasses \texttt{atexit}. The harness therefore
        re-emits its counters (distinct kernel PCs, comparison and
        return-token splices, record census) periodically on a line-buffered
        \texttt{stderr}, and the scheduler keeps the maximum.
\end{itemize}

\subsection{Measuring Coverage Without Double Counting}

The per-grain popcount of distinct kernel PCs, \emph{summed} over the fleet,
is biased: shared PCs (the dispatcher, session lookup, the VFS) are counted
once per grain, so the engine that cycles through the most grains wins
without reaching any new code. Each harness therefore keeps a $2^{17}$-bit
PC bitmap (index $= (\mathit{pc}\cdot c_3) \gg 47$, the same layout in both
tiers). Whenever the popcount grows, the bitmap is dumped, and
\texttt{gen.py} ORs every dump into a fleet-wide union. Its popcount,
$\mathit{kern\_pcs}_\cup$, is the primary metric of
Section~\ref{sec:eval}. \libfuzzer's $\mathit{ft}$ and corpus size are
reported only as secondary metrics: value coverage keeps more inputs by
construction, so those numbers favour the dataflow arms.

\subsection{Build and Target Bring-Up}

The kernel is built with \virtme from its default base plus an explicit list
of configuration items, so no machine-specific \texttt{.config} is needed:
\texttt{KCOV}, \texttt{KCOV\_ENABLE\_COMPARISONS},
\texttt{KCOV\_DATAFLOW\_\{ARGS,RET,INSTRUMENT\_ALL\}}, \texttt{SMB\_SERVER}
with SMBDirect and Kerberos, \texttt{UNICODE} (required for the casefold
path), Soft-RoCE and Soft-iWARP, \kasan (generic, inline), \texttt{SLUB\_DEBUG\_ON},
UBSAN bounds, lockdep with RCU proving, \texttt{DEBUG\_LIST}, hung-task and
soft-lockup panics, \texttt{FAILSLAB}, and DWARF5. \texttt{ksmbdzzer.py init}
then performs an idempotent nine-step bring-up: share and mount directories,
loading \ksmbd, provisioning the \texttt{fuzz} account, RDMA links,
\texttt{ksmbd.mountd}, the client mount, and an optional KDC. Grain harnesses
can be pre-compiled on the host (\texttt{build-grains}), which turns an 11--20
minute in-guest build over 9p into a cache hit.


\section{Evaluation}\label{sec:eval}

The evaluation answers five questions:
\begin{description}[nosep, font=\normalfont\bfseries]
  \item[RQ1] Do grains actually reach \ksmbd handler code?
  \item[RQ2] Does value-aware feedback reach more kernel code than PC-only
             feedback at the same cost per execution?
  \item[RQ3] Does input-to-state mutation from kernel comparison operands add
             coverage over havoc?
  \item[RQ4] Does \toolname find real vulnerabilities, and of which kinds?
  \item[RQ5] What do the instrumentation and the fleet cost?
\end{description}

\paragraph{What is and is not compared.} None of these questions is a
benchmark against another fuzzer. Edge-coverage engines and \toolname differ
in their unit of input (system call versus SMB procedure), in their feedback
(\emph{where} execution went versus \emph{what values} it carried), and in
their oracle (crash versus integrity effect). A difference in their numbers
could not be attributed to any one of these factors, and the bug classes that
\toolname targets, integrity violations that return success, do not register
as findings in a crash-only engine. RQ2 and RQ3 therefore ablate
\emph{within} \toolname: the arms share the kernel, the instrumentation, the
grains, the scheduler and the oracles, and differ only in whether kernel
values enter the coverage map and whether comparison operands drive mutation.
The \texttt{kcov} arm uses mainline \tracepc through the same harness, so it
places stock edge coverage on the same scale without changing anything else.
The goal is to isolate what knowledge of function inputs and outputs
contributes, not to rank fuzzers.

\subsection{Experimental Setup}

The kernel is linux-next 7.3.0-rc5 (\texttt{next-20261001}) with the \kcovdf
and \ksmbd remote-coverage patches, built with clang/LLVM~24.0.0git
carrying the \traceargs/\traceret pass and the configuration of
Section~\ref{sec:impl}. It runs under \virtme (QEMU/KVM) with
\pending{N}~vCPUs and 16\,GB of memory, on a \pending{host CPU} host.
Grains talk to \ksmbd over guest-internal loopback, so no network device is
involved.

\paragraph{Methodology.} Following the recommendations of Klees et
al.~\cite{klees2018evaluating} and Schloegel et al.~\cite{schloegel2024sok},
each arm runs \pending{$T$} independent trials. Every trial starts
\emph{cold}: the guest corpus, the host-durable mirror and the stability
baselines are wiped, and the trial gets a distinct but reproducible seed, so
no arm inherits another arm's inputs. The arms share the grain set, the
kernel image and the per-grain time cap of \pending{$G$}\,s. Phase~P3 is
disabled for the coverage comparison, because combination is a refinement
that would add a confound without contributing to the coverage question. I
report the median $\mathit{kern\_pcs}_\cup$ with its range, the two-sided
Mann--Whitney $U$ test, and the Vargha--Delaney
$\hat{A}_{12}$ effect size~\cite{arcuri2011practical}. The fleet is restricted to the 60
grains of \texttt{ALIGNED\_SUBSET.txt}, the grains tagged ALIGNED in at least
6 of the 9 runs of an earlier calibration campaign, so that a grain that
reaches the kernel only intermittently cannot bias one arm.

\subsection{Correctness Verification (RQ1)}

\texttt{selftest} runs each of the 210 grains four times with a
deterministic 64-byte seed per (grain, try) and takes the largest number of
distinct kernel PCs that the grain's own connection produced
(Table~\ref{tab:selftest}). On the recorded run, \textbf{208 grains reached
\ksmbd code and none was DEAD or BAIL}. The two exclusions are architectural:
\texttt{ipc} drives the generic-netlink channel to \texttt{ksmbd.mountd}
rather than an SMB connection, so the IP-keyed remote handle never fires, and
\texttt{rdma} needs an SMBDirect data plane that the loopback self-test does
not drive. Both stay in the fleet for real campaigns. Single-operation
grains that had intermittently looked DEAD reached 2{,}000--8{,}000 distinct
PCs once the asynchronous merge was awaited (artifact~(i) below).

\begin{table}[t]
\centering
\caption{Grain self-test verdicts, full 210-grain fleet (recorded
2026-07-22).}\label{tab:selftest}
\small
\begin{tabular}{@{}lrl@{}}
\toprule
\textbf{Verdict} & \textbf{Grains} & \textbf{Condition} \\
\midrule
WORKS & 208 & best-of-4 distinct kernel PCs $\geq 1$ \\
DEAD  & 0   & ran, but its connection produced no kernel PC \\
BAIL  & 0   & a prerequisite (pool, fid, auth) failed on every try \\
EXCL  & 2   & \texttt{ipc}, \texttt{rdma}: not measurable on a per-connection handle \\
\bottomrule
\end{tabular}
\end{table}

Three artifacts of running the whole fleet in one process initially made
correct grains look broken, and each has to be controlled for the
measurement to be accurate.
(i)~\ksmbd merges a remote section into the user buffer in
\texttt{kcov\_df\_remote\_stop()}, \emph{after} the response has been sent,
so a single-operation grain can read its buffer before the merge has landed.
The self-test therefore polls the buffer, idempotently, up to six times at
3\,ms intervals.
(ii)~Grains that tear down a connection (\texttt{encrypt},
\texttt{session\_setup}, \texttt{smb1\_*}, \texttt{logoff*}, \texttt{tdis*})
left the shared pool with a dead socket, and every pool grain after them then
bailed. Re-authenticating before every grain is not a fix, because it causes
the \texttt{mountd} storm described above. Instead, the pool is healed once,
only after a grain has bailed on every try.
(iii)~Per-connection coverage cannot measure \texttt{ipc} and \texttt{rdma};
hence EXCL rather than DEAD. Every grain that had looked DEAD or BAIL turned
to WORKS once these artifacts were controlled. No grain was removed to reach
this result.

\emph{Caveat.} The self-test was recorded on an earlier build of the same
tree, before the \kcovdf v4 UAPI. It must be repeated on the evaluation kernel:
\pending{re-run selftest on 7.3.0-rc5}.

\subsection{Information Gain Over Existing Tools (RQ2, RQ3)}

Table~\ref{tab:ablation} is the controlled ablation. Comparing
\texttt{dataflow-vec} with \texttt{pc-i2s} answers RQ2: both arms pay for the
same instrumentation and the same mutator, and differ only in whether values
are folded into coverage. Comparing \texttt{pc-i2s} with \texttt{pc-havoc}
answers RQ3. That comparison is meaningful only if input-to-state mutation was
actually active, so the table reports $\mathit{cmp\_hits}$, the number of
inputs in which a kernel comparison operand was spliced. If
$\mathit{cmp\_hits}=0$, the two arms are expected to tie, and a difference
would be noise. The \texttt{kcov} row places stock mainline \tracepc
coverage, routed through the same per-connection handle, on the same scale.

\begin{table}[t]
\centering
\caption{Engine ablation on the aligned subset: median fleet-union kernel PCs
over $T$ trials, with [min--max]. $p$ and $\hat{A}_{12}$ compare each arm with
\texttt{pc-i2s}. Corpus size and \texttt{max\_ft} are engine-biased secondary
metrics.}\label{tab:ablation}
\small
\begin{tabular}{@{}lrrrrrr@{}}
\toprule
\textbf{Arm} & $\mathbf{kern\_pcs_\cup}$ & $\mathbf{p}$ & $\hat{\mathbf{A}}_{12}$ & \textbf{cmp\_hits} & \textbf{corpus} & \textbf{max\_ft} \\
\midrule
\texttt{dataflow-vec} & \pending{TBD} & \pending{} & \pending{} & \pending{} & \pending{} & \pending{} \\
\texttt{dataflow-rel} & \pending{TBD} & \pending{} & \pending{} & \pending{} & \pending{} & \pending{} \\
\texttt{dataflow}     & \pending{TBD} & \pending{} & \pending{} & \pending{} & \pending{} & \pending{} \\
\texttt{pc-i2s}       & \pending{TBD} & ---        & ---        & \pending{} & \pending{} & \pending{} \\
\texttt{pc-havoc}     & \pending{TBD} & \pending{} & \pending{} & 0 (by design) & \pending{} & \pending{} \\
\texttt{kcov} (mainline) & \pending{TBD} & \pending{} & \pending{} & \pending{} & \pending{} & \pending{} \\
\bottomrule
\end{tabular}
\end{table}

\paragraph{Expected shape and how to read it.} Two properties of \ksmbd matter
when the table is read. First, return-value tokens are structurally rare in
this target. \ksmbd's C functions return NTSTATUS codes, error codes or
pointers, all of which the token filter discards, while file identifiers and
resume keys travel through \emph{output buffers}, and so appear in argument
records rather than in returns. A $\mathit{ret\_hits}$ near zero is therefore
expected, and it is not a failure of the mutator. Second, the value region is
bounded by design, so a dataflow arm can only win by reaching new
(call site, value class) pairs, not by inflating its feature count. The
$\mathit{kern\_pcs}_\cup$ metric counts PCs only, never values, so a value
arm is credited only for the \emph{code} that its value-driven corpus
reaches. \pending{Result paragraph: observed medians, significance, and the
grains that account for the difference.}

\subsection{Performance Overhead (RQ5)}

Two costs compound: the kernel instrumentation, and SMB round trips over
loopback. The foundational measurements apply unchanged. \kcovdf stores a
record in 18--21\,ns, and whole-kernel argument and return instrumentation
adds 66\% syscall latency with collection idle, on top of a \kasan
kernel~\cite{kim2026kcovdataflow}. Table~\ref{tab:perf} reports what this
means for a grain: the per-grain execution rate, the records per execution,
and the share of records that are comparisons, which is what drives the 2\,MiB
($2^{18}$-word) raw-tier buffer size.

\begin{table}[t]
\centering
\caption{Throughput per grain (median over the aligned subset).}\label{tab:perf}
\small
\begin{tabular}{@{}lrrr@{}}
\toprule
\textbf{Configuration} & \textbf{exec/s} & \textbf{records/exec} & \textbf{cmp share} \\
\midrule
\kcovdf, \texttt{INSTRUMENT\_ALL}, \texttt{dataflow-vec} & \pending{} & \pending{} & \pending{} \\
\kcovdf, \texttt{INSTRUMENT\_ALL}, \texttt{pc-i2s}      & \pending{} & \pending{} & \pending{} \\
\kcovdf, \ksmbd-only instrumentation                     & \pending{} & \pending{} & \pending{} \\
mainline KCOV (\texttt{kcov})                           & \pending{} & ---        & ---        \\
\bottomrule
\end{tabular}
\end{table}

The fast-bail threshold makes the cost of a misaligned grain explicit. A grain
that does not reach \ksmbd runs \emph{in-process} at more than 20{,}000
executions per second, while a grain that makes real SMB round trips stays far
below that rate. Before fast-bail, misaligned grains that ran until
\libfuzzer's execution cap accounted for most of the wall-clock time in a
dataflow-arm campaign while contributing almost no kernel PCs.
\pending{Report the measured fraction of wall-clock reclaimed by fast-bail.}

\subsection{Triage Capability}

Each finding carries the \kcovdf records of the execution that produced it.
For the contract oracle, the emitted reproducer contains the request bytes,
the NTSTATUS, the effect on the backing file system, and the decoded record
stream, including the arguments that \texttt{smb2\_create\_open\_flags()}
received for the probing \texttt{CREATE}. This makes it immediately clear
whether the SMB layer ever saw the dangerous access bits. For memory-safety
findings, the records just before the \kasan report tell the analyst which
length, count or index a function received, without adding \texttt{printk}
and without a crash dump. Values at function boundaries survive
optimization~\cite{kim2026kcovdataflow}. The fix for the
\texttt{ksmbd\_alloc\_user()} out-of-bounds read
(Section~\ref{sec:realworld}) explicitly credits \toolname and its
$\mathit{pc}\oplus\mathit{val}$ signal: \texttt{hash\_sz} is a 16-bit field
of an IPC response, and its out-of-range value is visible as a value even
where the edges it takes are not new.


\section{Case Study}\label{sec:realworld}

\subsection{Methodology}

All findings were reproduced on a sanitizer kernel, root-caused, and reported
to the \ksmbd and SMBDirect maintainers with a patch where possible.
Table~\ref{tab:bugs} lists the defects that \toolname campaigns surfaced and
that are now fixed in Linux.

\begin{table}[t]
\centering
\caption{Vulnerabilities found by \toolname and fixed upstream. CVE identifiers are CVE-2026-\textit{NNNNN}.}\label{tab:bugs}
\footnotesize
\setlength{\tabcolsep}{4pt}
\begin{tabular}{@{}llllr@{}}
\toprule
\textbf{Defect} & \textbf{Component} & \textbf{Class / detector} & \textbf{Fix} & \textbf{CVE (CVSS)} \\
\midrule
foreign-session \texttt{fp} put via wrong table & \texttt{vfs\_cache.c} & UAF / \kasan & \texttt{e7188199eff4} & 74522 (8.8) \\
\texttt{hash\_sz} not bounded by \texttt{hash[]}   & \texttt{user\_config.c} & slab OOB read / \kasan & \texttt{7405d0ba2943} & 90174 (7.1) \\
CQ freed with \texttt{ib\_destroy\_cq()}         & smbdirect \texttt{connection.c} & UAF in CQ poll work & \texttt{fe2c0cacbcff} & 90173 (9.8) \\
pools freed before QP on accept failure & smbdirect \texttt{accept.c} & leak + NULL deref & \texttt{383a9480f5f4} & 90172 (7.5) \\
child sockets released under listener lock        & smbdirect \texttt{socket.c} & UAF / \kasan + lockdep & \texttt{db82fbe4bb68} & 90171 \\
IPC fields read before length check               & \texttt{transport\_ipc.c} & slab OOB read / \kasan & \texttt{e9b33376bd07} & 90170 \\
\midrule
recursive \texttt{listen.lock} in cleanup         & smbdirect \texttt{socket.c} & lockdep & \texttt{76fa42c004eb} & --- \\
\bottomrule
\end{tabular}
\end{table}

\paragraph{File-table lifetime (CVE-2026-74522).} \texttt{ksmbd\_close\_fd()}
marked a file pointer closed and dropped its IDR reference, but it left the
pointer in the owning session's IDR. A \emph{different} session holding a
reference obtained through \texttt{ksmbd\_lookup\_fd\_inode()} could then
perform the final put through the wrong file table and free the object, while
the owner's IDR still pointed at it, and the owner's teardown then used freed
memory. The bug needs two sessions acting on one inode, which is the kind of
state that shared-object combination and the teardown oracle are designed to
create. The maintainer's fix removes the IDR entry under the owner's table lock.

\paragraph{IPC validation (CVE-2026-90174, CVE-2026-90170).} Both bugs are
violations of one contract: the kernel trusted lengths inside a
\texttt{ksmbd.mountd} response before checking that the response was large
enough to contain them. \texttt{ksmbd\_alloc\_user()} used the 16-bit
\texttt{hash\_sz} to size a \texttt{memcpy} from a fixed-size \texttt{hash[]}.
\texttt{ipc\_validate\_msg()} read \texttt{payload\_sz},
\texttt{session\_key\_len} and \texttt{ngroups} before checking
\texttt{msg\_sz} against the base structure of the event, so an 8-byte
response caused a read at offset~8 of a \texttt{kmalloc-8} object. The fixes
bound \texttt{hash\_sz} by \texttt{sizeof(hash)} and check the size of each
event's base structure before any dereference.

\paragraph{SMBDirect teardown (CVE-2026-90173, -90172, -90171).} The RDMA
transport produced three teardown-ordering bugs. Completion queues allocated
with \texttt{ib\_alloc\_cq\_any()} in \texttt{IB\_POLL\_WORKQUEUE} mode were
freed with \texttt{ib\_destroy\_cq()}, which does not cancel the poll work, so
a late completion after \texttt{rdma\_destroy\_qp()}, such as a Soft-RoCE RNR
error, re-queued work on a freed CQ. On accept failure, the memory pools and
the receive slab were destroyed while a posted receive was still owned by the
queue pair. Pending child sockets were released under the listener's handler
mutex and before the listener's \texttt{rdma\_destroy\_id()}, so
\texttt{\_cma\_cancel\_listens()} walked a freed child; the nested mutex also
produced a lockdep recursive-locking report. A separate fix
(\texttt{76fa42c004eb}) removes a recursive acquisition of
\texttt{listen.lock} during cleanup. All of these bugs lie on error and
teardown paths that are reached mainly under connection churn, which the
fleet's constant creation and destruction of sessions plausibly supplies as a
side effect.

\paragraph{Contract oracle.} \pending{Cross-user oracle outcome on the
evaluation kernel, with positive-control status.} If the oracle stays silent
while the positive control passes, the result is a \emph{negative} one with a
known meaning: for the three write-side procedures and the account mapping
tested, \ksmbd did not act with more authority than POSIX grants. The
SMB-layer checks themselves are a separate matter. \ksmbd honours lock, delete-on-close and overwrite requests on handles
whose \emph{granted access} lacks the corresponding right, contrary to
MS-SMB2. Those are real missing SMB-layer checks, but POSIX gates them, so they
cannot be used against a file the user does not own. I report them as
conformance defects, not as privilege escalation.

\subsection{Correction Note}\label{sec:correction}

Three corrections to earlier statements, two of them in the project's own
records, are part of the result.

\paragraph{Owner-on-own-file ``findings'' were false positives.} The first
contract oracle flagged a violation whenever \ksmbd accepted a dangerous
\texttt{CREATE} (for example \texttt{DELETE\_ON\_CLOSE} without
\texttt{DELETE} in the granted mask) on a victim file that the fuzzing account
had created itself. It produced five ``findings'', and all five were false
positives: \ksmbd authorizes at the VFS layer, and an owner may legitimately
truncate or delete its own file. The oracle was redesigned around a
root-owned victim and an effect-based verdict (Section~\ref{sec:oracles}).
The owner sweep is still run, but now only to drive coverage of
\texttt{smb2\_create\_open\_flags()} and input-to-state mapping, and it never
emits a finding.

\paragraph{Earlier comparison campaigns were invalid.} Every engine-comparison
summary recorded between 2026-08-28 and 2026-09-02 reports
$\mathit{kern\_pcs}_\cup = 0$ for all arms and was rejected by the runner's
all-zero guard. There were two causes. First, grains killed at saturation
never reached their \texttt{atexit} metric dump; periodic emission fixed
this. Second, the logs show
\texttt{KCOV\_DF\_REMOTE\_ENABLE(0x4b440001)} failing with \texttt{EFAULT}.
The \kcovdf v2/v4 UAPI takes the handle \emph{by pointer}, as a
\texttt{\_\_u64}, so that the full 64-bit handle survives 32-bit and compat
callers, and it moved the value count to header bits 32--47, while the harness
still passed the handle by value and parsed the count from bits 24--27. Those
campaigns therefore received no kernel coverage at all. The results in
Section~\ref{sec:eval} must come only from campaigns run after the harness
tracks the v4 UAPI. The same logs also show that the contract oracle never
ran in those campaigns: its probe failed to authenticate to
\texttt{[privtest]}, and the oracle disabled itself as designed. The lesson
generalizes: a fuzzer whose feedback channel
fails silently is a random-input generator that still looks like it is
working, and the all-zero guard is what exposed the failure here.

\paragraph{Oplock break deadlock.} A deadlock in level-II oplock breaking that
was observed during campaigns had already been fixed upstream, independently
(\texttt{27a9bc968d4a}: \ksmbd no longer holds \texttt{ci->m\_lock} while
waiting for a lease-break acknowledgement). I do not count it as a finding of
this work.


\section{Discussion \& Limitations}\label{sec:discussion}

\paragraph{Why not a head-to-head benchmark.} An obvious objection is that
\toolname is not compared with \syzkaller or with a network-protocol fuzzer.
The omission is deliberate. Such a comparison would mix a different input
model, a different feedback signal and a different oracle, so it would rank
two systems without explaining either. The question this work asks is
narrower and, I argue, more useful: does knowledge of the target's function
arguments and return values make mutation more effective than relying on
fuzzer-generated inputs alone? That question is answered by ablation with
everything else held fixed. The approach is also compatible with edge-coverage
engines rather than competing with them: \kcovdf records arrive through a
separate device and leave KCOV untouched~\cite{kim2026kcovdataflow}, so a
syscall-level fuzzer could adopt the same value fold and operand ring. Doing
so is a natural next step.

\paragraph{A direction worth developing.} Guiding mutation by what the kernel
actually received and returned is still in its early days, and this paper uses
it in a deliberately simple form: first-seen value classes, literal operand
splicing, boundary drives and return tokens. Several extensions follow
directly:
\begin{itemize}[nosep]
  \item \emph{Transform-aware correspondence}, matching an observed value
        against its UTF-16LE, big-endian or ASN.1-length encoding in the
        input.
  \item \emph{Distance feedback}, rewarding comparison operands that move
        closer rather than only operands that become equal.
  \item \emph{Field-level state}, using the struct fields \kcovdf already
        expands (for example lease state, credit counts or open flags) as a
        state model for sequencing procedures.
  \item \emph{Value-timed concurrency}, firing the second request of a race
        on the \emph{entry} record of a target function rather than after a
        fixed delay.
  \item \emph{Value-level contracts}, generalizing the integrity oracle to
        pre- and postconditions on the arguments and returns of security
        checks~\cite{kim2026kcovdataflow}.
\end{itemize}
Each of these uses information that an edge-coverage engine does not have.

\paragraph{Grains are hand-written.} Each grain is a few dozen lines of C, and
a fleet of 210 represents real effort. A declarative PDU specification
interpreted by one generic harness (\texttt{\$GRAIN\_SPEC}) already exists as
an opt-in path, and migrating the single-request grains to it would make
adding a procedure a one-line change. Stateful grains (compound, races,
durable reconnects, SPNEGO) still need code.

\paragraph{Fixed prefixes trade breadth for depth.} By construction, a grain
does not fuzz its own prefix, so a bug that requires a malformed negotiate
\emph{followed by} a deep operation is outside any single grain.
Pre-authentication surfaces are covered by grains whose target \emph{is} the
early PDU (\texttt{negotiate}, \texttt{session\_setup}, \texttt{spnego\_asn1}),
and combination covers two-procedure interactions. Sequences of three or more
procedures are not explored.

\paragraph{Input-to-state needs literal correspondence.} Operand splicing
fires when a wire field reaches a comparison or an argument unchanged, or as a
plain little-endian integer. SMB carries names in UTF-16LE, and NTLM and
SPNEGO wrap values in ASN.1 and HMACs. Values behind a MAC cannot be solved
this way at all.

\paragraph{Shallow dereference.} \kcovdf expands one level of struct fields.
\texttt{smb2\_open()} receives a \texttt{struct ksmbd\_work *} whose request
lives behind an opaque \texttt{void *}, so the request's fields become visible
only one frame deeper, where callees take them by value~\cite{kim2026kcovdataflow}.
For \ksmbd this costs little, because the handlers pass decoded fields down
promptly.

\paragraph{Throughput.} Every execution makes real SMB round trips over
loopback through a sanitizer kernel. Snapshot-based
designs~\cite{schumilo2021nyx,song2020agamotto,schumilo2022nyxnet} could
restore the post-prefix state directly and would avoid the per-execution
session cost. They would also remove the session churn that surfaced the
SMBDirect teardown bugs.

\paragraph{Concurrency.} The race mode of combination creates contention
without controlling interleavings. Razzer~\cite{jeong2019razzer} and
KRACE~\cite{xu2020krace} control interleavings, and value-timed races are
one of the extensions listed above.

\paragraph{Generality.} Nothing in the routing or the fold is specific to SMB.
Any in-kernel network service whose work runs off the client task (NFSD, the
TLS handshake upcall, in-kernel RPC) can adopt the
destination-address-keyed remote handle and the grain structure.

\paragraph{Ethics and disclosure.} All testing ran on loopback in disposable
VMs, against no system other than the author's own. Every finding in
Table~\ref{tab:bugs} was reported to the upstream maintainers, and each is
fixed in Linux.


\section{Related Work}\label{sec:related}

\paragraph{Kernel fuzzing.} \syzkaller~\cite{vyukov2015syzkaller} and
kAFL~\cite{schumilo2017kafl} use edge or block coverage. HEALER~\cite{sun2021healer},
SyzVegas~\cite{wang2021syzvegas}, Moonshine~\cite{pailoor2018moonshine} and
ACTOR~\cite{fleischer2023actor} improve sequence construction, scheduling,
seeds and action modelling. StateFuzz~\cite{zhao2022statefuzz} identifies
driver state variables statically and rewards new value combinations of
those variables. That is the closest idea to the value region here, but it
requires a static analysis to choose the variables, whereas \toolname rewards
values at every instrumented boundary and bounds the reward with value
classes. All of these systems work at the syscall interface. \toolname works
one layer higher, at the protocol procedure, and treats \syzkaller as
complementary rather than as a baseline.

\paragraph{Network-protocol fuzzing.} AFLNet~\cite{pham2020aflnet} uses
response codes as state, StateAFL~\cite{natella2022stateafl} infers state from
memory snapshots, SGFuzz~\cite{ba2022sgfuzz} tracks state variables, and
Nyx-Net~\cite{schumilo2022nyxnet} uses incremental snapshots. They target
user-space servers. The grain is a coarser but kernel-compatible answer to the
same problem: instead of learning to reach a state, each grain is written to
start in one.

\paragraph{Input-to-state and value steering.} RedQueen~\cite{aschermann2019redqueen},
AFL++ CmpLog~\cite{fioraldi2020aflpp} and Eclipser~\cite{choi2019eclipser}
solve comparisons cheaply from observed operands. Angora~\cite{chen2018angora}
and GREYONE~\cite{gan2020greyone} use taint. IJON~\cite{aschermann2020ijon}
lets an analyst annotate state. \toolname takes its operands from the
kernel's \tracecmp stream, adds a persistent operand ring and injection for
the case where the input does not hold the operand literally, and uses
function \emph{arguments}, not only comparisons, as steering targets.

\paragraph{Semantic and permission bugs.} Hydra~\cite{kim2019hydra} finds
semantic file-system bugs with checkers such as crash consistency.
PeX~\cite{zhang2019pex} and Crix~\cite{lu2019crix} find missing permission
checks statically. The contract oracle here is a dynamic, effect-based check
of a network file server's authorization contract. It reports a violation
only if the backing file system shows the effect, which avoids the false
positives that a status-based oracle produced (Section~\ref{sec:correction}).

\paragraph{Instrumentation.} KCOV~\cite{vyukov2016kcov} and \sancov~\cite{llvmsancov}
provide the coverage substrate, and \kcovdf~\cite{kim2026kcovdataflow}
provides the argument, return and comparison signal used here. This paper is
the first end-to-end use of that signal to fuzz a production kernel
subsystem.


\section{Conclusion}\label{sec:conclusion}

\toolname fuzzes the Linux in-kernel SMB server at the level of protocol
procedures. A grain starts every input past the protocol wall. A
per-connection remote handle delivers each worker exactly the kernel
arguments, returns and comparison operands that its own requests caused. A
bounded fold turns those values into coverage without destroying the
saturation signal, and the comparison operands drive input-to-state
mutation. Effect-based contract and stability oracles look for the bugs that
do not crash. 208 of 210 grains reach \ksmbd code. Campaigns found six CVEs in
file-table lifetime, IPC validation and SMBDirect teardown, all fixed
upstream. The work does not try to out-benchmark edge-coverage engines; it
explores a different axis, in which mutation is guided by what the target's
functions actually receive and return rather than by fuzzer-generated input
alone. A controlled multi-arm ablation, run on the same kernel and grains,
isolates that contribution. I believe the direction is promising and deserves
further development. Source code is available at \anonrepo{}.

%Bibliography
\bibliographystyle{unsrt}
\bibliography{../shared/references}

\appendix
\section{Dataflow-Directed Mutator}\label{app:i2s}

\begin{algorithm}[h]
\caption{\textsc{MutateI2S}: one directed step (library tier).}\label{alg:i2s}
\begin{algorithmic}[1]
\Function{MutateI2S}{$D$, $\mathit{seed}$}
  \If{arm $=$ \texttt{pc-havoc} \textbf{or} $\mathit{seed} \bmod 3 \neq 0$} \Return \textsc{Havoc}($D$) \EndIf
  \State grow $D$ to $\geq 16$ bytes (zero-filled, bounded by \texttt{max\_len})
  \State $U \gets \emptyset$; $h \gets 0$ \Comment{unplaced operands; hit count}
  \ForAll{record $r \in \mathit{buf}$ with $r.\tau = \texttt{C}$, width $w \geq 2$, $a_1 \neq a_2$}
    \If{$\mathrm{LE}_w(a_1)$ or $\mathrm{LE}_w(a_2)$ occurs in $D$}
      \State overwrite one occurrence with the other operand; $h \gets h+1$
    \Else
      \State $U \gets U \cup \{a_1, a_2\}$
    \EndIf
  \EndFor
  \For{$t \gets 1$ \textbf{to} $6$} \Comment{inject: this execution or persistent ring}
    \State $v \gets$ ($t$ odd \textbf{and} $U \neq \emptyset$) ? pick($U$) : pick($\mathit{CmpRing}$)
    \State write $\mathrm{LE}_{|v|}(v)$ at offset $(\mathit{seed} + 29t) \bmod (|D| - |v| + 1)$; $h \gets h+1$
  \EndFor
  \ForAll{record $r$ with $r.\tau = \texttt{E}$, field $v_k \notin \{0, \bot\}$}
    \For{$w \in \{8, 4, 2\}$ with $v_k < 2^{8w}$}
      \If{$\mathrm{LE}_w(v_k)$ occurs in $D$ at offset $o$}
        \State $x \gets$ pick from $\mathit{RetDict}$, this execution's returns, or $\mathit{Boundary}_w$
        \State write $\mathrm{LE}_w(x)$ at $o$; $h \gets h+1$; \textbf{break}
      \EndIf
    \EndFor
  \EndFor
  \State \Return $h > 0$ ? $D$ : \textsc{Havoc}($D$)
\EndFunction
\end{algorithmic}
\end{algorithm}

\section{Generated Library-Tier Grain Harness}\label{app:harness}

\begin{lstlisting}[language=C,caption={The harness \texttt{gen.py} emits for every library grain (abridged). Kernel coverage, not user-space coverage, populates \libfuzzer's extra counters.},label=lst:grain]
__attribute__((used, section("__libfuzzer_extra_counters")))
static uint8_t ctr[4096];
static int g_octet = 1;

int LLVMFuzzerInitialize(int *argc, char ***argv) {
    const char *ip = getenv("GRAIN_IP");          /* 127.0.0.<n> */
    if (ip) { const char *p = strrchr(ip, '.'); if (p) g_octet = atoi(p + 1); }
    if (pfz_init((unsigned long)g_octet) < 0)      /* mmap kcov_dataflow, publish */
        _exit(1);                                  /* handle 0x4b440000 | n        */
    return 0;
}

int LLVMFuzzerTestOneInput(const uint8_t *data, size_t size) {
    static uint32_t feats[4096];
    pfz_grain_run(IDX, data, (int)size);           /* fixed prefix + target PDU    */
    int n = pfz_get_features(feats, 4096);         /* fold this conn's records     */
    memset(ctr, 0, sizeof(ctr));
    for (int i = 0; i < n; i++) ctr[feats[i] & 4095]++;
    return 0;
}

size_t LLVMFuzzerCustomMutator(uint8_t *d, size_t s, size_t m, unsigned seed) {
    size_t r = pfz_mutate_i2s(d, s, m, seed);      /* Algorithm 1, else havoc      */
    return r ? r : LLVMFuzzerMutate(d, s, m);
}
\end{lstlisting}

\end{document}
